%
\newpage
\section{Experimental Setup}
\label{sec:experimental-setup}

\subsection{Protocol Simulation}
\begin{itemize}
    \item A single ``all-knowing'' orchestrator process (with access to all key material) simulates every party and performs the setup plus one full protocol round; no networking is involved.
    \item Round structure, parameterized over the number of users $\eta$ (\texttt{benchmark} directory):
    \begin{itemize}
        \item each virtual user holds a unique ID and a secret key share; the joint public key is produced by OpenFHE's chained multiparty key generation (each user extends the previous user's public key share; the final key is the joint public key);
        % FLAG: earlier notes said the public key shares are all "synthesized" into one key --- the generation is sequential/chained, not a combination of independent shares.
        \item each user submits an RLWE encryption of its value and $D$ one-hot RGSW index vectors under the joint key;
        % FLAG: earlier notes said each user "encodes their ID in an RGSW"; in the benchmark the written value is a plaintext value encrypted as RLWE, and the indices are the one-hot RGSW vectors of sec. impl:spar.
        \item the server executes \Call{Server::Write}{} once per user;
        \item decryption is multiparty: every user produces a partial decryption; fusion yields the plaintext.
    \end{itemize}
    \item The setup phase (key generation, state initialization) is considered free and excluded from all measurements.
\end{itemize}

\subsection{Unit Tests and Benchmarks}
\begin{itemize}
    \item Unit tests (\texttt{test} directory) empirically validate correctness: RGSW encryption and both products are tested for the plaintexts $0$, $1$ and a larger message, in both coefficient-packed and SIMD variants; the full protocol is tested parameterized over $\eta$ --- a passing run certifies the parameter set handles $\eta$ users.
    \item Tests check correctness at every stage and thereby localize failures (which operation failed and why), at the cost of speed.
    \item Benchmarks perform no correctness checks. One-shot timings suffice: client-side operations run independently per user, so the single-run time is the per-user latency, and for the server write only the per-user time is of interest.
    \item Tests and benchmarks share parameters and code paths, so a passing test validates exactly the computation the benchmark measures.
    \item The moduli $q_i$ are randomly generated per run; empirically the results do not depend on the specific $q_i$, only on the bit size of $Q$.
\end{itemize}

\subsection{Noise Measurement}
\begin{itemize}
    \item The testing environment always has access to the secret key, so RLWE noise is measured directly:
    \begin{equation*}
        x = (c_0, c_1) = (c_1 s + \Delta m + \epsilon,\; c_1)
        \implies \epsilon = c_0 - c_1 s - \Delta m, \qquad m = \Call{Decrypt}{x, s}.
    \end{equation*}
    % FLAG: state the reported metric explicitly (log2 of the infinity norm of eps, per limb?) --- ch. 4 tables depend on it.
    \item RGSW ciphertexts are intermediate objects; two ways to measure their noise:
    \begin{enumerate}
        \item extract the row encoding $\mu \cdot 1$ (exists whenever the gadget contains the entry $1$) and measure it as an RLWE ciphertext;
        \item compute the external product with a noiseless $\Call{RLWE}{1}$ (via \Call{MakePublicRGSW}{}-style zero-noise encryption) and measure the result.
    \end{enumerate}
    \item The latter is used: it applies to arbitrary gadgets, whereas the former assumes $1$ is a gadget entry --- traditional, but not guaranteed (and not true for the BV gadget rows scaled by $\tilde{e}_j$).
\end{itemize}
